% !TEX root = ../mainfile.tex
% !TEX encoding = UTF-8 Unicode
% !TEX spellcheck = en_US

\chapter{Introduction}
\label{ch:introduction}

Within a single year, two major shocks moved global commodity and financial markets. On the 2nd of April 2025, widely referred to as "Liberation Day", the United States Government announced broad-based tariffs on virtually all of its trading partners, citing persistent trade deficits and national security concerns. The policy included a base 10\% tariff and, subsequently, additional higher rates for trading partners with large surpluses against the U.S., such as China and the European Union. Following the announcement, the dollar fell sharply against a broad basket of currencies, with the Euro, Yen and Swiss Franc appreciating against it---the Franc by more than 3\% within the first week. The move was accompanied by a spike in equity-market volatility and a U.S. equities sell-off.

Less than a year later, on 28~February 2026, the United States and Israel launched airstrikes on Iran, beginning the 2026 Iran war. As escalation in the Gulf threatened to close the Strait of Hormuz and sent crude sharply higher, the dollar did exactly what the safe-haven literature predicts---it appreciated.

The puzzle speaks to a live debate about whether the dollar's exceptional status is eroding. A growing literature argues that the currency's safe-haven premium has weakened, pointing to the unusual co-movement of a falling dollar, rising Treasury yields and rising risk premia during the 2025 tariff episode \parencite{ostry_lloyd_corsetti_2025,kamin_2025,jiang_krishnamurthy_lustig_richmond_2026}. Against this, the long tradition on the dollar's ``exorbitant privilege'' and its role as the world's reserve and safe-haven currency \parencite{gourinchas_rey_2005,obstfeld_zhou_2023,habib_stracca_2012,ranaldo_soderlind_2010} implies the status should be durable. The two 2025--2026 shocks offer an unusually clean natural experiment for adjudicating the debate, because they differ in \emph{kind}: one is a U.S.\ trade-policy shock, the other a geopolitical oil-supply shock.

The central question of this thesis is therefore: \emph{does the dollar's safe-haven response depend on the type of shock---even when the United States is involved in both?} To adjudicate it, the design sets the two shocks against a historical benchmark---three episodes in all. The Liberation Day tariff announcement is a purely U.S.-originated economic-policy shock. The 2026 Israel--Iran/Hormuz crisis is a geopolitical oil-supply shock in which the United States is involved as a co-belligerent but still functions as a financial safe harbour. And the 2022 Russian invasion of Ukraine---an external military/oil shock not authored by the United States---serves as a control for the ``classical'' safe-haven response. Comparing the three isolates whether the dollar's behaviour is conditioned by the nature of the shock or merely by who caused it.

The thesis answers the question with two complementary empirical layers on a single daily dataset spanning January~2019 to June~2026. An \emph{event study} measures the cumulative abnormal response of the dollar, a cross-section of currency baskets, gold, equities and crude oil around each shock, and tests formally whether the dollar's response differs across events. A complementary \emph{regression} then tests whether the dollar's relationship with risk itself changed across the episodes---whether its safe-haven sensitivity to volatility held or broke. The two layers are complementary: the event study isolates the dollar's move around each discrete, dated shock, while the regression tests whether its underlying sensitivity to risk shifted across the episodes.

The contribution is threefold. First, to the best of the author's knowledge this is the first systematic side-by-side comparison of the Liberation Day and Iran/Hormuz episodes for currency markets. Second, the thesis bridges the safe-haven and commodity-currency literatures by carrying the oil channel and a gold benchmark through the same design, so that the dollar's behaviour can be read against gold and the oil-exposed currency cross-section rather than in isolation. Third, the rebuilt carry/dollar risk factors and the WMR currency panel extend standard cross-sectional tools through to mid-2026, allowing the most recent shocks to be studied with consistent, citable data.

The remainder of the thesis is organised as follows. Chapter~\ref{ch:literature} reviews the safe-haven, dollar-erosion and commodity-currency literatures. Chapter~\ref{ch:framework} sets out the conceptual framework---the working definition of a safe-haven currency, a classification that separates a shock's origin from its nature, and the hypotheses that follow. Chapter~\ref{ch:background} narrates the 2022--2026 events. Chapter~\ref{ch:methodology} describes the data and the event-study and regression methodologies. Chapter~\ref{ch:results} presents the empirical results. Chapter~\ref{ch:discussion} discusses the findings and their implications, and Chapter~\ref{ch:conclusion} concludes.

The results point to a dollar whose safe-haven status is conditional on the type of shock rather than on its originator. The two U.S.-involved shocks produce a mirror image: the Liberation Day tariffs depreciate the broad dollar index by over three and a half percent within twenty trading days, while the Hormuz oil-supply crisis appreciates it---by as much as the external Ukraine benchmark, from which its response is statistically indistinguishable. The complementary regression confirms the pattern: the dollar's risk sensitivity, significantly positive on normal days, collapses to essentially zero inside the tariff window yet holds through the Hormuz and Ukraine episodes. Gold emerges as a hedge against the dollar rather than a universal safe haven, and the oil channel sorts the currency cross-section only for the supply-driven shock.
